\documentclass[conference]{IEEEtran}
\IEEEoverridecommandlockouts

\usepackage{cite}
\usepackage{amsmath,amssymb,amsfonts}
\usepackage{graphicx}
\usepackage{textcomp}
\usepackage{xcolor}
\usepackage{booktabs}
\usepackage{multirow}
\usepackage{url}
\usepackage[hidelinks]{hyperref}
\usepackage{tabularx}
\usepackage{array}
\renewcommand{\topfraction}{0.9}
\renewcommand{\bottomfraction}{0.8}
\renewcommand{\textfraction}{0.07}
\renewcommand{\floatpagefraction}{0.7}

\newcommand{\todonote}[1]{\textcolor{red}{\textbf{[PENDIENTE: #1]}}}

\def\BibTeX{{\rm B\kern-.05em{\sc i\kern-.025em b}\kern-.08em
    T\kern-.1667em\lower.7ex\hbox{E}\kern-.125emX}}

\begin{document}

\title{Métodos constructivistas para agendamiento de salas de cirugía}

\author{
\IEEEauthorblockN{Juan Felipe Gallo\IEEEauthorrefmark{1}}
\IEEEauthorblockA{\IEEEauthorrefmark{1}Faculty of engineering, University of Antioquia\\
Medell\'in, Colombia \\
juanf.gallo@udea.edu.co}
}

\maketitle

\begin{abstract}
	LOREM
\end{abstract}


\begin{IEEEkeywords}
Metaheurísticas, Métodos constructivistas, GRASP, Medical Surgeon Schedule
\end{IEEEkeywords}

\section{Introduction}
\label{sec:introduction}
LOREM
\subsection{Research questions}

\subsection{Contributions}



The rest of this article is organized as follows: Section~\ref{sec:related_work}  summarizes related work; Section~\ref{sec:methodology} describes the data pipeline and the congruence
index methodology; Section~\ref{sec:results}  presents the results; Section~\ref{sec:discussion}  discusses limitations and implications; and Section~\ref{sec:conclusion} concludes.

\section{Related work}
\label{sec:related_work}

\section{Metodología}
\label{sec:methodology}

\subsection{Formulación del problema}

Sean $S$ el conjunto de cirugías, $R$ las salas, $K$ los cirujanos, $D$ los días y $B$ los bloques (AM y PM). Para cada cirugía $s$ se define el conjunto $F_s$ de opciones factibles $(d,b,r,k)$. Una opción es factible si cumple cuatro condiciones: el día $d$ está entre los días en que $s$ puede programarse, la sala $r$ está equipada para el tipo de $s$, el cirujano $k$ está certificado para ese tipo y el cirujano está disponible el día $d$. La decisión se representa con variables binarias $x_{s,d,b,r,k}$, definidas solo para $(d,b,r,k)\in F_s$:

\begin{equation}
	\sum_{(d,b,r,k)\in F_s} x_{s,d,b,r,k} \le 1 \quad \forall s
	\label{eq:cirugia}
\end{equation}
\begin{equation}
	\sum_{s}\sum_{k} x_{s,d,b,r,k} \le 1 \quad \forall r,d,b
	\label{eq:sala}
\end{equation}
\begin{equation}
	\sum_{s}\sum_{r} x_{s,d,b,r,k} \le 1 \quad \forall k,d,b
	\label{eq:cirujano}
\end{equation}

La restricción~\eqref{eq:cirugia} permite dejar una cirugía sin programar, \eqref{eq:sala} impide que una sala aloje dos cirugías en el mismo día y bloque, y \eqref{eq:cirujano} hace lo mismo con los cirujanos.

El objetivo es lexicográfico y sus criterios se comparan en este orden:

\begin{enumerate}
	\item $N$: número de cirugías programadas (se maximiza).
	\item $D$: dispersión, $D=\sum_{dep}\left(|Z_{dep}|-1\right)$, donde $Z_{dep}$ es el conjunto de zonas en que el departamento $dep$ tiene cirugías (se minimiza).
	\item $M$: número de cirugías fuera de la zona mayoritaria de su departamento (se minimiza).
\end{enumerate}

Las soluciones se comparan como la tupla $(N,-D,-M)$. El tercer criterio existe porque $D$ solo cambia cuando un departamento deja de usar una zona, por lo que un movimiento individual casi nunca la modifica. En cambio, $M$ varía con cada movimiento y le da a la búsqueda local una señal para avanzar.

\subsection{Heurística constructiva}

El constructivo asigna una cirugía por iteración:

\begin{enumerate}
	\item \textbf{Selección.} Se escoge la cirugía pendiente con \emph{menos opciones factibles} en ese momento, recalculadas dinámicamente según la ocupación actual. Las cirugías sin ninguna opción se descartan y quedan sin programar.
	\item \textbf{Asignación.} Para la cirugía elegida, cada opción $o=(d,b,r,k)$ recibe un costo
	\begin{equation}
		c(o)=\mathrm{dem}(r,d,b)+\mathrm{dem}(k,d,b)+\lambda\,\pi(o),
		\label{eq:costo}
	\end{equation}
	donde $\mathrm{dem}$ cuenta cuántas opciones de las cirugías pendientes usarían ese recurso en ese bloque, de modo que se evita consumir lo escaso que otras cirugías necesitan. El término $\pi(o)$ vale 1 si la sala está en una zona que el departamento del cirujano aún no usa (habiendo ya usado otra) y 0 en caso contrario, con $\lambda=0{,}5$.
	\item Se elige la opción de menor costo. Con $\alpha=0$ el método es determinista.
\end{enumerate}

\subsection{Búsqueda local (VND)}

Se usa un descenso de vecindario variable con primera mejora: cuando un vecindario encuentra una mejora, esta se aplica y la búsqueda reinicia en N1. El proceso termina cuando ningún vecindario mejora (con un límite de 30\,s). La Tabla~\ref{tab:vecindarios} resume los seis vecindarios.

\begin{table}[t]
	\caption{Vecindarios del VND}
	\label{tab:vecindarios}
	\centering
	\small
	\begin{tabular}{@{}l p{4.6cm} l@{}}
		\hline
		Vecindario & Movimiento & Mejora \\
		\hline
		N1 & Insertar una cirugía sin programar en un hueco libre & $N$ \\
		N2 & Cadena de expulsión de profundidad 1 & $N$ \\
		N3 & Cadena de expulsión de profundidad 2 & $N$ \\
		N4 & Sacar una cirugía programada e insertar dos sin programar & $N$ \\
		N5 & Reubicar una cirugía (día, bloque, sala, cirujano) & $(D,M)$ \\
		N6 & Intercambiar las salas de dos cirugías del mismo día y bloque & $(D,M)$ \\
		\hline
	\end{tabular}
\end{table}

En una cadena de expulsión, la cirugía ocupa un hueco ya tomado, se expulsan las cirugías que estorban (a lo sumo dos: la de la sala y la del cirujano) y todas deben reinsertarse de forma recursiva. Si alguna no puede reinsertarse, el movimiento se deshace mediante una bitácora de cambios. Como las cadenas exigen reinsertar a todas las cirugías expulsadas, nunca cambian \emph{cuáles} cirugías quedan por fuera; N4 cubre ese caso.

\subsection{GRASP}

GRASP repite el constructivo con aleatorización: entre las opciones de la cirugía elegida se forma la lista restringida de candidatos $\{o : c(o)\le c_{\min}+\alpha\,(c_{\max}-c_{\min})\}$ y se escoge una al azar. Cada solución se mejora con el VND. La iteración inicial es el constructivo voraz ($\alpha=0$), seguida de 20 iteraciones con $\alpha=0{,}3$. Se conserva la mejor solución según la tupla $(N,-D,-M)$.

\subsection{Búsqueda de vecindarios grandes (LNS)}

El LNS parte de la solución del VND y repite 200 veces un ciclo de destruir y reparar:

\begin{enumerate}
	\item \textbf{Destruir.} Se toma una cirugía sin programar (o una al azar si todas están programadas) y se retiran hasta $k=4$ cirugías que ocupan salas o cirujanos que esa cirugía podría usar.
	\item \textbf{Reparar.} Se completa la solución con el mismo constructivo ($\alpha=0{,}2$) y se aplica el VND.
	\item \textbf{Aceptación.} Se acepta la nueva solución si no empeora la tupla, lo que permite moverse por mesetas. Se guarda la mejor solución encontrada.
\end{enumerate}

\subsection{Modelo exacto de referencia}

El óptimo se calcula con un modelo entero resuelto con HiGHS, que incluye las restricciones~\eqref{eq:cirugia}--\eqref{eq:cirujano}, en dos etapas. La primera maximiza $N$. La segunda fija ese valor y minimiza $D$ mediante las variables auxiliares $z_{dep,zona}$ (el departamento usa la zona) y $u_{dep}$ (el departamento tiene cirugías), con $x_{s,d,b,r,k}\le z_{dep(k),zona(r)}$ y $D=\sum z-\sum u$. El límite de tiempo es de 120\,s por etapa. El modelo exacto cubre los dos primeros niveles del objetivo, no $M$.

\subsection{Diseño experimental}

\begin{itemize}
	\item \textbf{Instancias reales:} las 10 entregadas con el taller (18 a 34 cirugías).
	\item \textbf{Corridas:} GRASP y LNS con 5 semillas independientes por instancia (1 a 5).
	\item \textbf{Métricas:} $N$, $D$ y $M$; la brecha $100\,(N_{opt}-N_{mejor})/N_{opt}$; y el tiempo de cómputo de cada método.
	\item \textbf{Validación:} toda solución se verifica contra las restricciones del problema antes de reportarse.
	\item \textbf{Junta clínica (Punto~2):} para cada departamento se calcula la intersección de los días disponibles de sus cirujanos y, para cada día candidato, cuántas cirugías pierden opciones o quedan bloqueadas. Se reserva el día de menor costo, se retira de la disponibilidad de los cirujanos del departamento y se vuelve a ejecutar el método.
\end{itemize}
\section{Resultados}
\label{sec:results}


Los experimentos se realizaron sobre las 10 instancias del taller, con entre 18 y 34 cirugías cada una. GRASP y LNS se ejecutaron con 5 semillas independientes por instancia, y se reporta la mejor solución según la tupla $(N,-D,-M)$. El modelo exacto probó optimalidad en las 10 instancias.

\subsection{Calidad de las soluciones}

La Tabla~\ref{tab:calidad} presenta, para cada instancia y método, el número de cirugías programadas $N$ y la dispersión $D$ en la forma $N/D$.

\begin{table}[t]
	\caption{Cirugías programadas y dispersión ($N/D$) por instancia y método}
	\label{tab:calidad}
	\centering
	\footnotesize
	\setlength{\tabcolsep}{3.5pt}
	\begin{tabular}{@{}rrcccccc@{}}
		\hline
		Inst. & $|S|$ & Óptimo & Constr. & VND & GRASP & LNS \\
		\hline
		1  & 18 & 18/0 & 18/1 & 18/0 & 18/0 & 18/0 \\
		2  & 20 & 20/0 & 20/1 & 20/0 & 20/0 & 20/0 \\
		3  & 26 & 24/0 & 24/1 & 24/1 & 24/0 & 24/0 \\
		4  & 28 & 18/1 & 17/1 & 18/1 & 18/1 & 18/1 \\
		5  & 30 & 23/0 & 21/1 & 23/1 & 23/0 & 23/0 \\
		6  & 30 & 14/0 & 14/0 & 14/0 & 14/0 & 14/0 \\
		7  & 30 & 14/1 & 14/2 & 14/2 & 14/1 & 14/1 \\
		8  & 32 & 12/3 & 12/3 & 12/3 & 12/3 & 12/3 \\
		9  & 30 & 15/0 & 15/1 & 15/1 & 15/0 & 15/0 \\
		10 & 34 & 20/1 & 20/3 & 20/2 & 20/1 & 20/1 \\
		\hline
		Total & 278 & 178/6 & 175/14 & 178/11 & 178/6 & 178/6 \\
		\hline
	\end{tabular}
\end{table}

\paragraph*{Ocupación.} El óptimo programa 178 de las 278 cirugías (64\,\%). Las instancias 1 y 2 se programan por completo y la 3 al 92\,\%, mientras que en las instancias 6 a 10 solo se programa entre el 38\,\% y el 59\,\%, por lo que en ellas los recursos, y no el método, limitan el resultado.

\paragraph*{GRASP y LNS.} Ambos alcanzaron el óptimo lexicográfico en las 10 instancias: brecha de 0\,\% en $N$ y la misma dispersión mínima que el modelo exacto ($D=6$ en total). En ambos métodos, las 5 semillas alcanzaron el mismo valor de $N$ en cada instancia.

\paragraph*{Constructivo.} Alcanzó el $N$ óptimo en 8 de las 10 instancias; en la instancia~4 programó una cirugía menos que el óptimo y en la~5, dos menos. Su dispersión total fue de 14, frente a 6 del óptimo, y solo en las instancias 6 y 8 coincidió con el óptimo lexicográfico.

\paragraph*{VND.} La búsqueda local cerró la diferencia en $N$ del constructivo y alcanzó el óptimo en las 10 instancias. En dispersión, en cambio, redujo el total de 14 a 11 sin llegar al óptimo, y coincidió con el óptimo lexicográfico en 5 de las 10 instancias (1, 2, 4, 6 y 8). Esto sugiere que los vecindarios N1 a N4, orientados a $N$, bastan en estas instancias, mientras que N5 y N6 dejan al VND atrapado en óptimos locales de dispersión que GRASP y LNS sí superan.

\subsection{Tiempos de cómputo}

La Tabla~\ref{tab:tiempos} resume los tiempos por método. Para GRASP y LNS se promedian las 5 semillas de cada instancia. El tiempo del VND corresponde a la búsqueda local partiendo del constructivo, y el del LNS, a la ejecución partiendo de la solución del VND. El tiempo de HiGHS incluye las dos etapas del modelo exacto.

\begin{table}[t]
	\caption{Tiempo de cómputo por método (s), sobre las 10 instancias}
	\label{tab:tiempos}
	\centering
	\small
	\begin{tabular}{@{}lccc@{}}
		\hline
		Método & Media & Mínimo & Máximo \\
		\hline
		Constructivo            & 0,006 & 0,001 & 0,021 \\
		VND (desde constructivo)& 0,017 & 0,001 & 0,087 \\
		GRASP                   & 0,131 & 0,055 & 0,295 \\
		LNS (desde VND)         & 1,886 & 0,352 & 7,484 \\
		HiGHS (dos etapas)      & 0,097 & 0,015 & 0,423 \\
		\hline
	\end{tabular}
\end{table}

El constructivo y el VND resolvieron todas las instancias en menos de 0,1\,s. GRASP tardó en promedio 0,13\,s. El LNS fue el más lento (1,9\,s en promedio), pero este promedio está dominado por las instancias 1 y 2 (7,5\,s y 6,4\,s); en las ocho restantes, el promedio es de 0,62\,s. Ambas instancias son aquellas en las que se programan todas las cirugías: GRASP se detiene en cuanto programa todas con dispersión cero, mientras que el LNS ejecuta siempre sus 200 iteraciones.

El modelo exacto fue más rápido que el LNS en las 10 instancias y más rápido que GRASP en 8; en las instancias 1 y 2 fue donde HiGHS tardó más (0,42\,s y 0,40\,s) y GRASP lo superó. En ningún caso HiGHS superó 0,5\,s.

\subsection{Junta clínica}

Para cada departamento se calculó el día común de menor costo (primero, menos cirugías bloqueadas; luego, menos cirugías afectadas), se retiró ese día de la disponibilidad de sus cirujanos y se volvió a resolver la instancia. Como ejemplo, en la instancia~1 reservar un día afecta entre 5 y 7 cirugías para Dep0, entre 3 y 5 para Dep1 y entre 14 y 18 para Dep2, y en ningún caso una cirugía queda sin opciones. Si los cirujanos de un departamento no comparten ningún día disponible, no se le reserva junta.

\begin{table}[t]
	\caption{Efecto de reservar la junta clínica de cada departamento}
	\label{tab:junta}
	\centering
	\small
	\begin{tabular}{@{}rcccc@{}}
		\hline
		Inst. & $N$ sin junta & $N$ con junta & Pérdida & $D$ (sin/con) \\
		\hline
		1  & 18 & 18 & 0 & 0/0 \\
		2  & 20 & 20 & 0 & 0/0 \\
		3  & 24 & 24 & 0 & 0/0 \\
		4  & 18 & 18 & 0 & 1/1 \\
		5  & 23 & 20 & 3 & 0/0 \\
		6  & 14 & 12 & 2 & 0/0 \\
		7  & 14 & 10 & 4 & 1/1 \\
		8  & 12 & 11 & 1 & 3/3 \\
		9  & 15 & 7  & 8 & 0/0 \\
		10 & 20 & 16 & 4 & 1/2 \\
		\hline
		Total & 178 & 156 & 22 & --- \\
		\hline
	\end{tabular}
\end{table}

Reservar la junta no tuvo costo en 4 de las 10 instancias (1 a 4). En total, las cirugías programadas pasan de 178 a 156, una pérdida de 22 (12,4\,\%). La pérdida es mayor en la instancia~9, donde se programan 7 cirugías en lugar de 15, y en las instancias 7 y 10, con 4 cirugías menos cada una. La dispersión se mantiene igual salvo en la instancia~10, donde sube de 1 a 2.

Estos valores deben leerse como una cota superior del costo real de la junta: el día de cada departamento se elige de forma independiente con una regla voraz, y la solución con junta proviene de una sola semilla de GRASP y de LNS, sin validación con el modelo exacto.
\section{Discussion}
\label{sec:discussion}

\section{Conclusiones}
\label{sec:conclusion}


Se abordó la programación semanal de quirófanos con un objetivo lexicográfico que prioriza el número de cirugías programadas y, después, la concentración de cada departamento en pocas zonas. Se desarrollaron una heurística constructiva, una búsqueda local de vecindario variable, GRASP y una búsqueda de vecindarios grandes, y se evaluaron frente al óptimo de un modelo entero en las 10 instancias reales del taller.

GRASP y LNS alcanzaron el óptimo lexicográfico en las 10 instancias. La búsqueda local (VND) alcanzó el número óptimo de cirugías en las 10, pero su dispersión total fue de 11 frente a 6 del óptimo, y el constructivo solo quedó por debajo del óptimo en $N$ en dos instancias. En estas instancias, entonces, la diferencia entre métodos está en la dispersión y no en el número de cirugías.

En tiempo de cómputo, el constructivo y el VND tardaron menos de 0,1\,s, GRASP 0,13\,s en promedio y LNS 1,9\,s, mientras que HiGHS resolvió todas las instancias en menos de 0,5\,s. Por tanto, con instancias de este tamaño las heurísticas no ofrecen una ventaja de tiempo sobre el método exacto, y estos experimentos no permiten concluir cómo se comportarían con instancias mayores.

Reservar un día común para la junta clínica de cada departamento no tuvo costo en 4 de las 10 instancias y redujo de 178 a 156 el total de cirugías programadas (12,4\,\%) con la estrategia propuesta. Este valor es una cota superior del costo real, porque el día de cada departamento se elige de forma voraz e independiente y la solución con junta no se validó con el modelo exacto.

Las conclusiones están limitadas por el tamaño de las instancias (entre 18 y 34 cirugías) y por el uso de una única configuración de parámetros para GRASP y LNS. Como trabajo futuro se propone: evaluar los métodos en instancias de mayor tamaño; incorporar el día de junta como decisión del modelo y de la búsqueda, y no como un paso previo; añadir un criterio de parada al LNS; e incluir el tercer criterio de desempate en el modelo exacto.


\bibliographystyle{IEEEtran}
\bibliography{references}

\end{document}
